\section{Forcing over Laurent series fields}\label{app:forcing}

Throughout this appendix $K$ satisfies~\textup{(L)}; from Section~\ref{frc:ssec:trace} on, $K$ is dense in $\kappa((X))$,
so $\Kh=\kappa((X))$ and $\abs\vartheta=\varrho$. The results of Sections~\ref{sec:sharpness} and~\ref{sec:classification}
proved from Section~\ref{frc:ssec:trace} on, except Lemma~\ref{obs:lem:single-vector}, which assumes this density, depend
only on $K$ as a normed field; Lemma~\ref{obs:lem:normalization} extends them to every $K$ satisfying~\textup{(L)}.
Theorem~\ref{cls:thm:spreading}(3) and the last assertion of Example~\ref{cls:ex:garling}(3) are proved in
Section~\ref{frc:ssec:spreading} over every nontrivially normed field of characteristic~$2$. Lemma~\ref{fib:lem:coeff}
holds, with the same proof, for every field $\kappa$ of characteristic $p$ in place of the finite field
of~\textup{(L)}; \cite{II} uses it in that generality.

\subsection{Coefficients and costs}\label{frc:ssec:costs}
For a normed space $E$ with $c_{00}(\N,K)\subseteq E\subseteq\linf(\N,K)$ and $\norm x_\infty\le\norm x$, the determinant
array $\Omega$ of \S\ref{ssec:ext}, with $S=\N$, satisfies $\abs{\Omega_z(c)}\le\norm z$, as in the proof of
Lemma~\ref{lem:proj}, so it extends to $\widehat\Lambda^k_\pi E$ with values in $\linf(\N^k,\kappa((X)))$. Likewise
$z(t,t')$ denotes the continuous extension to $E\widehat\otimes_\pi E$ of the entry $x(t)y(t')$ of $x\otimes y$.

For a matrix $\mathsf Z$ over $\kappa$ with rows indexed by a finite set $R$ and columns by a finite set $R'$,
$\operatorname{cost}(\mathsf Z)$ is the least value of
$\sum_i\abs{\operatorname{supp}g_i}^{1/2}\abs{\operatorname{supp}h_i}^{1/2}$ over the representations
$\mathsf Z=\sum_ig_ih_i^{\mathsf T}$ with $g_i\in\kappa^R$ and $h_i\in\kappa^{R'}$; it is subadditive.

\begin{lemma}[tournament lemma]\label{fib:lem:tournament}
Let $\mathsf Z$ be an $m\times m$ matrix over $\kappa$ such that every off-diagonal entry of $\mathsf Z-\mathsf Z^{\mathsf T}$ is nonzero. Then
$\operatorname{cost}(\mathsf Z)\ge m\ln m/(1+\ln(2\abs\kappa))$.
\end{lemma}

\begin{proof}
\emph{Step 1 (parts).} Let $\mathsf Z=\sum_ig_ih_i^{\mathsf T}$ with all $g_i,h_i$ nonzero, let $\mathcal V_i$ be the
union of the supports of $g_i$ and $h_i$, and for $t\in\mathcal V_i$ let the part $\mathcal V_i(t)$ be the set of points
of $\mathcal V_i$ at which $(g_i,h_i)$ is a multiple of $(g_i(t),h_i(t))$. The $(t,t')$ entry
$g_i(t)h_i(t')-h_i(t)g_i(t')$ of $g_ih_i^{\mathsf T}-h_ig_i^{\mathsf T}$ is nonzero exactly when $t,t'\in\mathcal V_i$ lie
in different parts; so every pair $t\ne t'$ lies in some $\mathcal V_i$, in different parts.

\emph{Step 2 (deletion).} Following Hansel's deletion argument~\cite{Ha} with weighted choices, choose independently
for each $i$ one part of $\mathcal V_i$, with probability proportional to its size, and delete the points of $\mathcal V_i$
outside it. At most one point survives, since two survivors would lie in different parts of some $\mathcal V_i$. The
point $t$ survives with probability $\exp(-L_t)$, where
$L_t=\sum_{i:t\in\mathcal V_i}\ln(\abs{\mathcal V_i}/\abs{\mathcal V_i(t)})$; so $\sum_t\exp(-L_t)\le1$, and convexity of
the exponential gives $\sum_tL_t\ge m\ln m$. Moreover $\sum_tL_t=\sum_i\operatorname{ent}_i$ with
$\operatorname{ent}_i=\sum_{t\in\mathcal V_i}\ln(\abs{\mathcal V_i}/\abs{\mathcal V_i(t)})$.

\emph{Step 3 (one term).} Exchanging $g_i$ and $h_i$, which does not change the parts, we may assume that the sizes
$N_1$ and $N_2$ of their supports satisfy $N_1\le N_2$, so $N=\abs{\mathcal V_i}\le2N_2$. It remains to show
$\operatorname{ent}_i\le(1+\ln(2\abs\kappa))\sqrt{N_1N_2}$. The points of $\mathcal V_i$ outside the support of $g_i$ form one
part, which contributes $(N-N_1)\ln(N/(N-N_1))\le N_1$. The other parts divide the support of $g_i$ by the value of
$h_i/g_i$ into at most $\abs\kappa$ sets; since a distribution on at most $\abs\kappa$ points has entropy at most
$\ln\abs\kappa$, they contribute at most $N_1\ln(N/N_1)+N_1\ln\abs\kappa$. Hence
\[
 \operatorname{ent}_i\le N_1\Bigl(1+\ln\frac{2\abs\kappa N_2}{N_1}\Bigr)\le(1+\ln(2\abs\kappa))\sqrt{N_1N_2},
\]
the second inequality with $x=N_2/N_1$, since $x^{-1/2}(1+\ln(2\abs\kappa x))$ is decreasing for $x\ge1$, as
$\ln(2\abs\kappa)>1$.
\end{proof}

\begin{lemma}[constant coefficients]\label{fib:lem:coeff}
Let $R,R'$ be finite sets and $r\in\{2,\infty\}$, let $z$ belong to the completion of
$\ell^r(R,K)\otimes_\pi\ell^r(R',K)$, and let $\coeff_0z\in\kappa^{R\times R'}$ be the matrix of constant coefficients of
its entries.
\begin{enumerate}
\item If $r=2$, then $\operatorname{cost}(\coeff_0z)\le\norm z/(1-\varrho^2)$.
\item If $r=\infty$, then $\rank(\coeff_0z)\le M_\varrho\norm z$.
\end{enumerate}
\end{lemma}

\begin{proof}
For $x\in\ell^r(R,K)$ and $n\in\Z$ let $x_{[n]}\in\kappa^R$ have entries $\coeff_n(x(t))$. Then
$\coeff_0(x\otimes y)=\sum_nx_{[n]}y_{[-n]}^{\mathsf T}$, a finite sum.
(1) The vector $x_{[n]}$ is supported in $O_n(x)=\{t:x(t)\ne0,\ \ord x(t)\le n\}$, and
$\sum_n\abs{O_n(x)}\varrho^{2n}=\norm x^2/(1-\varrho^2)$. By the Cauchy--Schwarz inequality,
$\operatorname{cost}\coeff_0(x\otimes y)\le\sum_n\abs{O_n(x)}^{1/2}\abs{O_{-n}(y)}^{1/2}\le\norm x\norm y/(1-\varrho^2)$.
(2) If $\norm x=\varrho^\nu$ and $\norm y=\varrho^{\nu'}$, only the $l+1$ values $\nu\le n\le-\nu'$ contribute, where
$l=-\nu-\nu'$, and none if $l<0$. Since $l+1\le M_\varrho\varrho^{-l}$, this gives
$\rank\coeff_0(x\otimes y)\le M_\varrho\norm x\norm y$.

For general $z$, write $z=\sum_jx_j\otimes y_j$ as an absolutely convergent series with
$\sum_j\norm{x_j}\norm{y_j}$ as close to $\norm z$ as we like, by approximating $z$ quickly with algebraic tensors and decomposing the
first approximation and the successive differences nearly optimally. The terms with $\norm{x_j}\norm{y_j}<1$ have
entries of absolute value at most $\varrho$, since $\abs{x(t)}\le\norm x$ and absolute values lie in $\varrho^\Z$; by
the nonarchimedean inequality so does their sum, which therefore does not change $\coeff_0$. The finitely many
remaining terms give the bounds with any number larger than $\norm z$ in place of $\norm z$, by subadditivity of cost
and rank, and hence with $\norm z$ itself.
\end{proof}

\begin{lemma}[two-slot slice]\label{fib:lem:slice}
Let $E$ be a normed space with $c_{00}(\N,K)\subseteq E\subseteq\linf(\N,K)$ and $\norm x_\infty\le\norm x$, let
$\gamma_1\ne\gamma_2$ be slots in $\{1,\dots,k\}$, and fix $c_\gamma\in\N$ for the remaining slots $\gamma$. There is a
bounded linear $\Upsilon\colon\widehat\Lambda^k_\pi E\to E\widehat\otimes_\pi E$ of norm at most $k(k-1)$ such that
$(\Upsilon z)(t,t')=\Omega_z(c)$, where $c$ has $t$ in slot $\gamma_1$, $t'$ in slot $\gamma_2$ and $c_\gamma$ in each slot
$\gamma$.
\end{lemma}

\begin{proof}
For $x_1,\dots,x_k\in E$, expansion along the rows $\gamma_1,\gamma_2$ gives
\[
 \det(x_b(c_a))_{a,b}=\sum_{b\ne b'}x_b(c_{\gamma_1})x_{b'}(c_{\gamma_2})\det_{bb'}(x),
\]
where $\det_{bb'}(x)$, the sum of
$\sgn(\sigma)\prod_{a\ne\gamma_1,\gamma_2}x_{\sigma(a)}(c_a)$ over the permutations with $\sigma(\gamma_1)=b$ and
$\sigma(\gamma_2)=b'$, does not involve $c_{\gamma_1},c_{\gamma_2}$, and
$\abs{\det_{bb'}(x)}\le\prod_{b''\ne b,b'}\norm{x_{b''}}$ by the nonarchimedean inequality. So the $k$\nobreakdash-linear
map $x\mapsto\sum_{b\ne b'}\det_{bb'}(x)\,x_b\otimes x_{b'}$ into $E\otimes_\pi E$ has norm at most $k(k-1)$, and its
values have the stated determinants as entries. It is alternating, since these entries vanish when two $x_b$ coincide
and a tensor in $E\otimes E$ with all entries zero is zero. Proposition~\ref{prop:exterior} gives $\Upsilon$ on
$\Lambda^k_\pi E$, and continuity extends it and the entry formula to the completion.
\end{proof}

\subsection{The \texorpdfstring{$\ell^r$}{l\textasciicircum r} threshold}\label{frc:ssec:threshold}
\begin{proof}[Proof of Theorem~\ref{fib:thm:ell2}]
\emph{Step 1 (a coefficient form).} For finitely supported $a_1,\dots,a_k,x_1,\dots,x_k$ with entries in $\F_p$ and
$c\in\N^k$ put
\[
 \Xi(a;x)(c)=\coeff_0\bigl(\Omega_{Q(a_1,\dots,a_k)(x_1\wedge\dots\wedge x_k)}(c)\bigr)\in\kappa,
\]
which is $\F_p$\nobreakdash-multilinear in $(a,x)$ and alternating in $x$ and in $c$. For such $b_1,\dots,b_k,x_1,\dots,x_k$
and $\lambda\in K^k$, with $a_\lambda=\sum_i\lambda_ib_i$, the identity
$Q(a_\lambda,\dots,a_\lambda)(x_1\wedge\dots\wedge x_k)=(a_\lambda x_1)\wedge\dots\wedge(a_\lambda x_k)$ is an equality of
polynomials in $\lambda$ with coefficients in $\widehat\Lambda^k_\pi E$. Comparing the coefficients of
$\lambda_1\cdots\lambda_k$ (Lemma~\ref{lem:idprinc}) and applying $\coeff_0\circ\Omega$ gives
\begin{equation}\label{fib:eq:pol1}
 \sum_{\sigma\in S_k}\Xi(b_{\sigma(1)},\dots,b_{\sigma(k)};x)(c)
 =\sum_{\sigma\in S_k}\det\bigl((b_{\sigma(j)}x_j)(c_i)\bigr)_{i,j},
\end{equation}
since these determinants lie in $\F_p$. The proofs of Lemmas~\ref{lem:H}, \ref{lem:C} and~\ref{lem:D} use only
multilinearity, alternation and~\eqref{fib:eq:pol1} of the map $\Psi$ of Appendix~\ref{app:FF}, and the
expansion~\eqref{eq:expand}. For $\Xi$ the coefficient function
$T(d;y;c)=\Xi(e_{d_1},\dots,e_{d_k};e_{y_1},\dots,e_{y_k})(c)$ takes values in the finite field $\kappa$
and~\eqref{eq:expand} follows from $\F_p$\nobreakdash-multilinearity, the inputs there having entries in $\{0,\pm1\}$; so
those three lemmas hold for $\Xi$. Since $k!=0$ in $\kappa$, Lemma~\ref{lem:D} shows that the cluster value $\chi$ is
not constant; as adjacent rank exchanges generate $S_k$, there are $\tau$ and $\ell$ with
$\delta=\chi(\tau)-\chi(\tau')\ne0$, where $\tau'=(\ell\ \ell{+}1)\circ\tau$ as in Lemma~\ref{lem:S}. Put
$\gamma_1=\tau^{-1}(\ell)$ and $\gamma_2=\tau^{-1}(\ell+1)$.

\emph{Step 2 (the quantitative bound).} Fix $m\ge1$, and let $a=(a_1,\dots,a_k)$ and $x=(x_1,\dots,x_k)$ be the inputs in
the proof of Lemma~\ref{lem:S}, with $\gamma_1,\gamma_2$ as the slots $\alpha,\beta$ there. By Lemma~\ref{lem:C}, the
$m\times m$ matrix $\mathsf Z$ of the values $\Xi(a;x)(c)$ with $A_n(1)$ in slot $\gamma_1$, $B_{n'}(1)$ in slot
$\gamma_2$ and $C_\gamma(1)$ in each other slot $\gamma$ is $\mathsf Z=\chi(\tau')\mathbf 1\mathbf 1^{\mathsf T}+\delta U$.
All entries of $a$ and $x$ lie in $\{0,\pm1\}$, and $x_{\gamma_1}$, $x_{\gamma_2}$ have $4m$ nonzero entries and the
other $x_\gamma$ four, so $\zeta=Q(a_1,\dots,a_k)(x_1\wedge\dots\wedge x_k)$ has $\norm\zeta\le\norm Q(4m)^{2/r}4^{(k-2)/r}$.
The slice $\Upsilon$ of Lemma~\ref{fib:lem:slice} for the slots $\gamma_1,\gamma_2$ and $c_\gamma=C_\gamma(1)$, followed by
restriction of both tensor factors to the coordinates $\{A_n(1)\}$ and $\{B_{n'}(1)\}$, gives $z$ with
$\norm z\le k(k-1)\norm\zeta$ and $\coeff_0z=\mathsf Z$. If $r=2$, Lemma~\ref{fib:lem:coeff}(1) gives
$\operatorname{cost}(\mathsf Z)\le k(k-1)2^km\norm Q/(1-\varrho^2)$, while $\mathsf Z-\mathsf Z^{\mathsf T}=\delta(U-U^{\mathsf T})$
and Lemma~\ref{fib:lem:tournament} give $\operatorname{cost}(\mathsf Z)\ge m\ln m/(1+\ln(2\abs\kappa))$. If $r>2$, the
$\ell^r$ norms dominate the supremum norms, and Lemma~\ref{fib:lem:coeff}(2) gives
$m-1\le\rank\mathsf Z\le M_\varrho k(k-1)\norm\zeta$, the first inequality since $\delta U$ is invertible. This
proves~\eqref{fib:eq:quant}, which fails for large $m$.

\emph{Step 3 (lifts and splittings).} A lift $Q_0$ as in Lemma~\ref{fib:lem:test}, for Banach values, would give such a
$Q$, namely $Q(a_1,\dots,a_k)=Q_0(\mu_{a_1},\dots,\mu_{a_k})$, since $\norm{\mu_a}\le\norm a_\infty$.
Lemma~\ref{fib:lem:test}, Theorem~\ref{thm:duality}(1) and the dichotomy of Proposition~\ref{prop:homogeneous-lift} give
the assertion about $A^k_{E,E;F}$, Proposition~\ref{fib:prop:split} the failure of the splittings, and
Corollary~\ref{cor:discrete-targets} the last assertion, since an equivalent nonarchimedean norm on
$\widehat\Lambda^k_\pi E$ would make $A^k_{E,E;\widehat\Lambda^k_\pi E}$ analytic.
\end{proof}

\begin{proof}[Proof of Corollary~\ref{fib:cor:threshold}]
Theorem~\ref{fib:thm:lower}, with the examples following it, and Theorem~\ref{fib:thm:ell2}.
\end{proof}

\begin{proof}[Proof of Corollary~\ref{fib:cor:pairs}]
The positive cases are Theorem~\ref{fib:thm:lower}. In the negative cases, repeat the proof of
Theorem~\ref{fib:thm:ell2} with $Q(a_1,\dots,a_k)=Q_0(\mu_{a_1},\dots,\mu_{a_k})$ for a lift $Q_0$ on $\cL(E,D)$ as in
Lemma~\ref{fib:lem:test}, with $\widehat\Lambda^k_\pi D$ in place of $\widehat\Lambda^k_\pi E$ and the slice formed in
$D\widehat\otimes_\pi D$. If $r_1\le r_2$, then $\norm{\mu_a}\le\norm a_\infty$ and
$\norm\zeta\le\norm{Q_0}(4m)^{2/r_1}4^{(k-2)/r_1}$. If $r_1>r_2$, H\"older's inequality gives
$\norm{\mu_a}\le\norm a_{r_0}$ with $1/r_0=1/r_2-1/r_1$; as $a_{\gamma_1}$, $a_{\gamma_2}$ have $2m$ nonzero entries and
the other $a_\gamma$ two, $\norm\zeta\le\norm{Q_0}(2m)^{2/r_0}2^{(k-2)/r_0}(4m)^{2/r_1}4^{(k-2)/r_1}$. In both cases
$\norm\zeta\le4^k\norm{Q_0}m^{2/\min(r_1,r_2)}$. If $r_2=2\le r_1$, Lemmas~\ref{fib:lem:coeff}(1)
and~\ref{fib:lem:tournament} give $m\ln m/(1+\ln(2\abs\kappa))\le k(k-1)4^km\norm{Q_0}/(1-\varrho^2)$; if $r_1,r_2>2$,
Lemma~\ref{fib:lem:coeff}(2) gives $m-1\le M_\varrho k(k-1)4^k\norm{Q_0}m^{2/\min(r_1,r_2)}$. Both fail for large $m$.
\end{proof}

\subsection{A trace norm over \texorpdfstring{$\kappa((X))$}{a Laurent series field}}\label{frc:ssec:trace}
Let $\mathcal O=\kappa[[X]]$, the closed unit ball of $\kappa((X))$. For finite sets $R,R'$ we identify
$z\in\kappa((X))^{R\times R'}$ with $\sum_{t,t'}z(t,t')\,e_t\otimes e_{t'}$; its projective norm $\norm z_\pi$ over
$\kappa((X))$ for the supremum norms on $\kappa((X))^R$ and $\kappa((X))^{R'}$ is the infimum of
$\sum_a\norm{g_a}_\infty\norm{h_a}_\infty$ over the finite representations $z=\sum_ag_ah_a^{\mathsf T}$. Fix enumerations $t_1,t_2,\dots$ of $R$ and $t'_1,t'_2,\dots$ of $R'$. Since $\mathcal O$ is
a principal ideal domain whose nonzero elements are units times powers of $X$, the elementary divisor theorem, applied to
$X^cz\in\mathcal O^{R\times R'}$ for a large integer $c$, writes every $z$ as
\begin{equation}\label{frc:eq:smith}
 z=U_1\Delta U_2,\qquad U_1\in GL_R(\mathcal O),\quad U_2\in GL_{R'}(\mathcal O),
\end{equation}
where $\Delta$ has the entries $X^{d_1},\dots,X^{d_{\rank z}}$, with integers $d_1\le\dots\le d_{\rank z}$, at the positions
$(t_1,t'_1),\dots,(t_{\rank z},t'_{\rank z})$ and zeros elsewhere. For $l\in\Z$ let $\rank_l(z)$ be the number of indices
$i$ with $d_i\le l$.

\begin{lemma}[a trace norm]\label{frc:lem:trace-norm}
Let $R,R'$ be finite sets and $z\in\kappa((X))^{R\times R'}$.
\begin{enumerate}
\item The numbers $\rank_l(z)$ do not depend on the factorization~\eqref{frc:eq:smith}.
\item \textup{(Level formula.)} $\norm z_\pi=\sum_i\varrho^{d_i}=(1-\varrho)\sum_{l\in\Z}\varrho^l\rank_l(z)$.
\item \textup{(Submatrices.)} If $z_0$ is the submatrix of $z$ with rows in $S\subseteq R$ and columns in
$S'\subseteq R'$, then $\rank_l(z_0)\le\rank_l(z)$ for every $l\in\Z$.
\item \textup{(Rank at a level.)} If $\abs R=\abs{R'}$, $e\in\Z$, and $z=X^ew$ for a matrix $w\in\mathcal O^{R\times R'}$
whose reduction modulo $X$ is invertible over $\kappa$, then $\rank_e(z)=\abs R$.
\end{enumerate}
\end{lemma}

\begin{proof}
For $l\ge0$ let $\mathcal O_l=\mathcal O/X^{l+1}\mathcal O$, a local ring with residue field $\kappa$, and let
$\operatorname{gen}(Y)$ be the least number of generators of a finitely generated $\mathcal O_l$\nobreakdash-module $Y$.
(a) If an $\mathcal O_l$\nobreakdash-module is generated by $N$ elements, so is each submodule: the module is a quotient of
$\mathcal O^N$, and the preimage of the submodule is free of rank at most $N$, since $\mathcal O$ is a principal ideal
domain. (b) $\operatorname{gen}(Y)=\dim_\kappa Y/XY$ by Nakayama's lemma; so
$\bigoplus_{i=1}^N\mathcal O_l/X^{c_i}\mathcal O_l$ needs exactly $N$ generators when every $c_i\ge1$.

\emph{Step 1 (generators).} Let $w\in\mathcal O^{R\times R'}$, $l\ge0$, and let $\operatorname{Col}_l(w)\subseteq
\mathcal O_l^R$ be the submodule generated by the columns of $w$ reduced modulo $X^{l+1}$. If $w=U_1\Delta U_2$ is a
factorization~\eqref{frc:eq:smith}, then $\Delta=U_1^{-1}wU_2^{-1}$ has entries in $\mathcal O$, so all $d_i\ge0$; we
claim $\operatorname{gen}(\operatorname{Col}_l(w))=\rank_l(w)$. The columns of $w$ and of $U_1\Delta$ generate the same
submodule, since $U_2$ is invertible over $\mathcal O$; multiplication by $U_1$ is an automorphism of $\mathcal O_l^R$; and
the columns of $\Delta$ generate a module isomorphic to $\bigoplus_{d_i\le l}\mathcal O_l/X^{l+1-d_i}\mathcal O_l$, which
needs exactly $\rank_l(w)$ generators by~(b).

\emph{Step 2 (item~(1)).} Choose $c\in\Z$ with $X^cz\in\mathcal O^{R\times R'}$. Multiplying~\eqref{frc:eq:smith} by $X^c$
gives a factorization of $X^cz$ with the exponents $d_i+c\ge0$. So $\rank_l(z)$ is $0$ if $l+c<0$ and
$\operatorname{gen}(\operatorname{Col}_{l+c}(X^cz))$ by Step~1 otherwise; neither depends on the factorization.

\emph{Step 3 (item~(2)).} The two right sides agree, since $(1-\varrho)\sum_{l\ge d}\varrho^l=\varrho^d$. Writing $z$ as
the sum over $i$ of $X^{d_i}$ times the product of the column $t_i$ of $U_1$ and the row $t'_i$ of $U_2$, vectors over
$\mathcal O$, gives a representation of cost at most $\sum_i\varrho^{d_i}$. Conversely, let $z\ne0$ and
$z=\sum_{a=1}^Ng_ah_a^{\mathsf T}$ with all $g_a,h_a$ nonzero. Since $\abs{\kappa((X))^\times}=\varrho^\Z$,
$g_a=X^{\nu_a}g'_a$ and $h_a=X^{\nu'_a}h'_a$ with vectors $g'_a,h'_a$ over $\mathcal O$ of supremum norm $1$; with
$m_a=\nu_a+\nu'_a$ the cost is $\sum_a\varrho^{m_a}$. Multiplying $z$ and every $g_a$ by $X^c$ multiplies the cost and
$\sum_i\varrho^{d_i}$ by $\varrho^c$, so we may assume every $m_a\ge0$; then $z\in\mathcal O^{R\times R'}$. For $l\ge0$,
$z\equiv\sum_{m_a\le l}X^{m_a}g'_ah_a'^{\mathsf T}$ modulo $X^{l+1}$, so $\operatorname{Col}_l(z)$ lies in the submodule
generated by the reductions of the $g'_a$ with $m_a\le l$, and (a) with Step~1 gives
$\rank_l(z)\le N_l=\abs{\{a:m_a\le l\}}$; for $l<0$ both are $0$. Hence
$\sum_a\varrho^{m_a}=(1-\varrho)\sum_l\varrho^lN_l\ge(1-\varrho)\sum_l\varrho^l\rank_l(z)$, and the infimum over
representations gives the lower bound.

\emph{Step 4 (item~(3)).} Multiplying $z$ and $z_0$ by the same $X^c$ replaces $l$ by $l+c$ on both sides, so we may
assume $z\in\mathcal O^{R\times R'}$ and $l\ge0$, both sides vanishing for $l<0$. Deleting columns replaces
$\operatorname{Col}_l(z)$ by a submodule, which needs no more generators by~(a); deleting rows replaces it by its image
under the coordinate projection $\mathcal O_l^R\to\mathcal O_l^S$, which is generated by the images of generators. Step~1 gives $\rank_l(z_0)\le\rank_l(z)$.

\emph{Step 5 (item~(4)).} Let $\Pi_0\in\mathcal O^{R\times R'}$ have the entry $1$ at each position $(t_i,t'_i)$ and $0$
elsewhere. The determinant of $\Pi_0^{-1}w$ reduces to a nonzero element of $\kappa$, so it is a unit and
$\Pi_0^{-1}w\in GL_{R'}(\mathcal O)$. Then $z=I\cdot(X^e\Pi_0)\cdot\Pi_0^{-1}w$ is a factorization~\eqref{frc:eq:smith}
with all $\abs R$ exponents equal to $e$, so $\rank_e(z)=\abs R$.
\end{proof}

\subsection{Coefficient functions of multiplier lifts}\label{frc:ssec:coefficients}
Throughout this subsection $G$ and $D$ are sequence spaces (Definition~\ref{obs:def:sequence-spaces}) with
$\norm{e_t}_G\le1$ and $\norm y_\infty\le\norm y_D$, $\Omega$ is the determinant array on $\widehat\Lambda^k_\pi D$ of
Section~\ref{frc:ssec:costs}, and $\widetilde Q\colon c_{00}(\N,K)^k\to\cL(\Lambda^k_\pi G,\widehat\Lambda^k_\pi D)$ is a
\emph{multiplier lift}: a $k$\nobreakdash-linear map with
\begin{equation}\label{frc:eq:multiplier-lift}
 \widetilde Q(b,\dots,b)(x_1\wedge\dots\wedge x_k)=(bx_1)\wedge\dots\wedge(bx_k)\qquad(b,x_1,\dots,x_k\in c_{00}(\N,K))
\end{equation}
and $\norm{\widetilde Q(e_{a_1},\dots,e_{a_k})}\le\varrho^{-l_0}$ for all $a\in\N^k$ and some integer $l_0\ge0$. For
$l\in\Z$ and $a,y,c\in\N^k$ put
\[
 \Omega_{a;y}(c)=\Omega_{\widetilde Q(e_{a_1},\dots,e_{a_k})(e_{y_1}\wedge\dots\wedge e_{y_k})}(c),\qquad
 T_l(a;y;c)=\coeff_l\Omega_{a;y}(c)\in\kappa .
\]
Since $\abs{\Omega_z(c)}\le\norm z$, $\norm{e_{y_1}\wedge\dots\wedge e_{y_k}}\le1$ and the nonzero absolute values of
$\kappa((X))$ lie in $\varrho^\Z$, a nonzero $\Omega_{a;y}(c)$ has order at least $-l_0$, so $T_l=0$ for $l<-l_0$. For
$g_1,\dots,g_k,h_1,\dots,h_k\in c_{00}(\N,K)$ and $c\in\N^k$, multilinearity gives the finite expansion
\begin{equation}\label{frc:eq:expand}
 \Omega_{\widetilde Q(g_1,\dots,g_k)(h_1\wedge\dots\wedge h_k)}(c)=\sum_{a,y\in\N^k}\prod_{\gamma=1}^kg_\gamma(a_\gamma)
 h_\gamma(y_\gamma)\cdot\Omega_{a;y}(c).
\end{equation}

\begin{lemma}[homogeneity]\label{frc:lem:homogeneity}
For every integer $l_1\ge0$ there is an infinite set $I\subseteq\N$ such that for every $l\le l_1$,
$T_l(\mathbf z)=T_l(\mathbf z')$ whenever $\mathbf z,\mathbf z'\in I^{3k}$ have the same pattern in the sense of Appendix~\ref{app:FF}.
\end{lemma}

\begin{proof}
The map $\mathbf z\mapsto(T_l(\mathbf z))_{-l_0\le l\le l_1}$ takes values in $\kappa^{l_0+l_1+1}$, a finite set whose
size is a power of $p$. Identifying it with a finite field of that size, Lemma~\ref{lem:H} with $N=3k$ gives $I$; for
$l<-l_0$, $T_l=0$.
\end{proof}

For the rest of this subsection fix an integer $l_1\ge0$ and $I$ as in Lemma~\ref{frc:lem:homogeneity}. As in
Appendix~\ref{app:FF}, clusters are the sets $C=\{C(1)<C(2)<C(3)<C(4)\}$ of four consecutive elements of $I$, disjoint
clusters are compared by $C<C'$, and the operator and vector weights $o=(1,0,-1,0)$ and $w=(1,-1,1,-1)$
satisfy~\eqref{eq:W}, with $o$ in place of $s$. Put $o_C=\sum_io_ie_{C(i)}$ and $w_C=\sum_iw_ie_{C(i)}$, so that
$o_Cw_C=e_{C(1)}-e_{C(3)}$. For $l\le l_1$ and pairwise disjoint clusters $C_1,\dots,C_k$ put
\begin{align*}
 \chi_l(C_1,\dots,C_k)&=\coeff_l\Omega_{\widetilde Q(o_{C_1},\dots,o_{C_k})(w_{C_1}\wedge\dots\wedge w_{C_k})}
 \bigl(C_1(1),\dots,C_k(1)\bigr)\\
 &=\sum_{i,i'\in[4]^k}\prod_{\gamma=1}^ko_{i_\gamma}w_{i'_\gamma}\cdot
 T_l\bigl(C_1(i_1),\dots,C_k(i_k);C_1(i'_1),\dots,C_k(i'_k);C_1(1),\dots,C_k(1)\bigr),
\end{align*}
by~\eqref{frc:eq:expand}, the weights being constants in $\{0,\pm1\}$. As in Appendix~\ref{app:FF}, the pattern of each
$3k$\nobreakdash-tuple in the sum is determined by $(i,i')$ and the order pattern $\tau\in S_k$ of $(C_1,\dots,C_k)$, so
by Lemma~\ref{frc:lem:homogeneity} the value depends only on $\tau$; we write it $\chi_l(\tau)$. Also $\chi_l=0$ for
$l<-l_0$.

\begin{lemma}[two-level cancellation]\label{frc:lem:clusters}
Let $\cC_1,\dots,\cC_k$ be finite families of clusters that are pairwise disjoint as sets of clusters, such that all
clusters in $\bigcup_\gamma\cC_\gamma$ are pairwise disjoint, and let $\mathrm{lev}^o_\gamma,\mathrm{lev}^w_\gamma\colon
\cC_\gamma\to\Z_{\ge0}$ be functions. Put
\[
 g_\gamma=\sum_{C\in\cC_\gamma}\vartheta^{\mathrm{lev}^o_\gamma(C)}o_C,\qquad
 h_\gamma=\sum_{C\in\cC_\gamma}\vartheta^{\mathrm{lev}^w_\gamma(C)}w_C\qquad(\gamma=1,\dots,k).
\]
Let $C_\gamma\in\cC_\gamma$ for $\gamma=1,\dots,k$, let $c=(C_1(1),\dots,C_k(1))$, and let
$\lambda=\sum_\gamma\bigl(\mathrm{lev}^o_\gamma(C_\gamma)+\mathrm{lev}^w_\gamma(C_\gamma)\bigr)$. Then for every integer
$l\le l_1$,
\[
 \coeff_l\Omega_{\widetilde Q(g_1,\dots,g_k)(h_1\wedge\dots\wedge h_k)}(c)
 =\sum_{l'\ge0}\coeff_{l'}(\vartheta^\lambda)\,\chi_{l-l'}(C_1,\dots,C_k),
\]
a finite sum.
\end{lemma}

\begin{proof}
\emph{Step 1 (expansion).} Expand by~\eqref{frc:eq:expand}. As the clusters are disjoint and $g_\gamma$, $h_\gamma$ are
supported in the union of the clusters of $\cC_\gamma$, a nonzero term has, in each slot $\gamma$, an operator index
$a_\gamma=\Cop_\gamma(i_\gamma)$ and a vector index $y_\gamma=\Cvec_\gamma(i'_\gamma)$ with
$\Cop_\gamma,\Cvec_\gamma\in\cC_\gamma$, and the coefficient $\prod_\gamma o_{i_\gamma}w_{i'_\gamma}\cdot\vartheta^{\lambda'}$,
where $\lambda'=\sum_\gamma\bigl(\mathrm{lev}^o_\gamma(\Cop_\gamma)+\mathrm{lev}^w_\gamma(\Cvec_\gamma)\bigr)\ge0$ depends
only on the \emph{cluster assignment} $(\Cop_\gamma,\Cvec_\gamma)_\gamma$. Since $\vartheta^{\lambda'}$ has nonnegative
order and $T_{l''}=0$ for $l''<-l_0$,
$\coeff_l(\vartheta^{\lambda'}\Omega_{a;y}(c))=\sum_{l'\ge0}\coeff_{l'}(\vartheta^{\lambda'})\,T_{l-l'}(a;y;c)$, a finite
sum.

\emph{Step 2 (output-free clusters).} Group the terms by cluster assignment and by $l'$. Suppose an assignment has an
occupied cluster $C'$ that is none of $C_1,\dots,C_k$. Then $C'$ belongs to exactly one family $\cC_\gamma$, so only
$a_\gamma$ and $y_\gamma$ can lie in $C'$, and no output index does. Fix $l'$ and all positions outside $C'$, and sum over
the positions inside $C'$. The factor $\coeff_{l'}(\vartheta^{\lambda'})$ and the weights of the other slots are common
factors, and $T_{l-l'}$ is pattern-homogeneous on $I$ since $l-l'\le l_1$. As $C'$ consists of consecutive elements of $I$,
cases (i)--(iii) of the proof of Lemma~\ref{lem:C} apply with $T_{l-l'}$ in place of $T$, and~\eqref{eq:W} makes the
partial sum vanish.

\emph{Step 3 (the surviving assignment).} In a surviving assignment every occupied cluster is an output cluster, and the
only output cluster in $\cC_\gamma$ is $C_\gamma$; so $\Cop_\gamma=\Cvec_\gamma=C_\gamma$ for every $\gamma$,
$\lambda'=\lambda$, and the contribution is $\sum_{l'\ge0}\coeff_{l'}(\vartheta^\lambda)\chi_{l-l'}(C_1,\dots,C_k)$.
\end{proof}

\begin{lemma}[the diagonal]\label{frc:lem:diagonal}
For every integer $l\le l_1$, $\sum_{\tau\in S_k}\chi_l(\tau)$ equals $1$ if $l=0$ and $0$ otherwise.
\end{lemma}

\begin{proof}
\emph{Step 1 (polarization).} Let $b_1,\dots,b_k,x_1,\dots,x_k\in c_{00}(\N,K)$ and $a_\alpha=\sum_i\alpha_ib_i$ for
$\alpha\in K^k$. Both sides of~\eqref{frc:eq:multiplier-lift} with $b=a_\alpha$ are polynomials in $\alpha$, and
comparing the coefficients of $\alpha_1\cdots\alpha_k$ by Lemma~\ref{lem:idprinc} gives
\begin{equation}\label{frc:eq:polarized}
 \sum_{\sigma\in S_k}\widetilde Q(b_{\sigma(1)},\dots,b_{\sigma(k)})(x_1\wedge\dots\wedge x_k)
 =\sum_{\sigma\in S_k}(b_{\sigma(1)}x_1)\wedge\dots\wedge(b_{\sigma(k)}x_k).
\end{equation}

\emph{Step 2 (right side).} Take clusters $C_1<\dots<C_k$, $b_\gamma=o_{C_\gamma}$, $x_\gamma=w_{C_\gamma}$ and
$c=(C_1(1),\dots,C_k(1))$, and apply $\Omega_\cdot(c)$ to~\eqref{frc:eq:polarized}. On the right, a permutation
$\sigma\ne\id$ has a slot $\gamma$ with $\sigma(\gamma)\ne\gamma$, and $o_{C_{\sigma(\gamma)}}w_{C_\gamma}=0$ by
disjointness; the identity contributes $\det\bigl((o_{C_b}w_{C_b})(C_a(1))\bigr)_{a,b}=1$, since
$o_Cw_C=e_{C(1)}-e_{C(3)}$.

\emph{Step 3 (left side).} On the left, the sign manipulation in the proof of Lemma~\ref{lem:D}, which uses only that the
wedge and $\Omega_z$ are alternating, identifies the summand of $\sigma$ with
\[
 \Omega_{\widetilde Q(o_{C_{\sigma(1)}},\dots,o_{C_{\sigma(k)}})(w_{C_{\sigma(1)}}\wedge\dots\wedge w_{C_{\sigma(k)}})}
 \bigl(C_{\sigma(1)}(1),\dots,C_{\sigma(k)}(1)\bigr),
\]
whose $l$th coefficient is $\chi_l(\sigma)$, since $(C_{\sigma(1)},\dots,C_{\sigma(k)})$ has order pattern $\sigma$.
Taking $\coeff_l$ gives $\sum_\sigma\chi_l(\sigma)=\coeff_l(1)$.
\end{proof}

\begin{corollary}[an adjacent transposition]\label{frc:cor:transposition}
There are $\tau\in S_k$ and $\ell\in\{1,\dots,k-1\}$ with $\chi_0(\tau)\ne\chi_0(\tau')$, where $\tau'=(\ell\ \ell{+}1)\circ\tau$.
For such $\tau,\ell$ put $\gamma_1=\tau^{-1}(\ell)$, $\gamma_2=\tau^{-1}(\ell+1)$, $\delta_l=\chi_l(\tau)-\chi_l(\tau')$ for
$l\le l_1$, and let $l_*$ be the least $l\le l_1$ with $\delta_l\ne0$. Then $-l_0\le l_*\le0$.
\end{corollary}

\begin{proof}
Otherwise $\chi_0$ is invariant under the adjacent transpositions, which generate $S_k$, hence constant, and
Lemma~\ref{frc:lem:diagonal} with $l=0$ gives $k!\,\chi_0=1$, impossible since $k!=0$ in $\kappa$. As $\delta_0\ne0$ and
$\delta_l=0$ for $l<-l_0$, $-l_0\le l_*\le0$.
\end{proof}

\subsection{The product triangle bound}\label{frc:ssec:ptb}
For a symmetric sequence space $D$ and finite $J\subseteq\N$, $D_J$ is the span of the $e_t$, $t\in J$; we use three facts.
\begin{enumerate}
\item[(P1)] For finite $J,J'\subseteq\N$, a tensor in $D_J\otimes D_{J'}$ has the same projective norm in
$D_J\otimes_\pi D_{J'}$ as in $D\otimes_\pi D$, since multiplication by $1_J$ is a contractive projection of $D$ onto $D_J$
(Lemma~\ref{obs:lem:multiplier}(1)).
\item[(P2)] Relabeling the rows and the columns of finitely supported tensors by injective maps $\N\to\N$, and
multiplying rows and columns by sequences in $K$ of absolute value $1$, are isometries for the norm of $D\otimes_\pi D$:
by symmetry of the norm they map each $D_J$ isometrically onto some $D_{J''}$, and (P1) applies.
\item[(P3)] For finite $J,J'\subseteq\N$, the completion of $D_J\otimes_\pi D_{J'}$ is $\Kh^{J\times J'}$, with a norm
between the maximum and the sum of the absolute values of the entries, since $\norm\cdot_\infty\le\norm\cdot_D$ and
$\norm{e_t}_D=1$; as $K$ is dense in $\Kh$ and this norm is continuous for the maximum norm, $\norm{\lambda z}=\abs\lambda\norm z$
for $\lambda\in\Kh$. Multiplication by $1_J\otimes1_{J'}$ extends to a contraction of $D\widehat\otimes_\pi D$ onto this
completion whose value at $z$ is the restriction of $z$ to $J\times J'$ (Section~\ref{frc:ssec:costs}).
\end{enumerate}

\begin{proof}[Proof of Theorem~\ref{obs:thm:ptb}]
Let $C_Q=\norm Q$, $l_0=\max(0,\lceil\log_{1/\varrho}C_Q\rceil)$ and $C_\star=64\cdot8^{k-2}k(k-1)\,C_Q+\varrho^{-l_0}+1$.
We show \textup{(PTB)} with constant $C_\star$. For general $K$ satisfying~\textup{(L)},
Lemma~\ref{obs:lem:normalization} then gives \textup{(PTB)} with $\abs\vartheta$ in place of $\varrho$ in $C_\star$ and $l_0$.

\emph{Step 0 (a multiplier lift).} For $a_1,\dots,a_k\in c_{00}(\N,K)$ put
$\widetilde Q(a_1,\dots,a_k)=Q(\mu_{a_1},\dots,\mu_{a_k})$, a $k$\nobreakdash-linear map with
\begin{equation}\label{frc:eq:mult-bound}
 \norm{\widetilde Q(a_1,\dots,a_k)}\le C_Q\prod_{\gamma=1}^k\norm{\mu_{a_\gamma}}_{E\to D}.
\end{equation}
Since $Q(u,\dots,u)=\Lambda^k_\pi u$, $\widetilde Q$ satisfies~\eqref{frc:eq:multiplier-lift}; since $\norm{e_t}_E=1$,
$\norm y_\infty\le\norm y_D$ and $\norm{\mu_{e_t}}_{E\to D}=1$ (Lemma~\ref{obs:lem:multiplier}(1),(3)), it is a multiplier
lift with $G=E$, as $\norm{\widetilde Q(e_{a_1},\dots,e_{a_k})}\le C_Q\le\varrho^{-l_0}$. The rest of the proof uses $Q$
only through these properties of $\widetilde Q$ and~\eqref{frc:eq:mult-bound}.

\emph{Step 1 (configuration).} Assume first:
\begin{enumerate}
\item[(A)] $\operatorname{supp}a=\operatorname{supp}x=J_1$ and $\operatorname{supp}b=\operatorname{supp}y=J_2$ with
$J_1\cap J_2=\emptyset$, and every nonzero entry of $a$, $x$, $b$ and $y$ is a nonnegative power of $\vartheta$.
\end{enumerate}
If one of the four vectors is $0$, both sides of~\eqref{obs:eq:ptb} vanish; so let them be nonzero and put
$N_0=\norm{\mu_a}_{E\to D}\norm x_E\norm{\mu_b}_{E\to D}\norm y_E>0$. List $J_1\cup J_2$ as $f_1<\dots<f_m$, and let
$l_1\ge0$ be the least integer with $m^2\varrho^{l_1+1}\le N_0$. Take $I$ as in Lemma~\ref{frc:lem:homogeneity} for this
$l_1$, and $\tau$, $\ell$, $\gamma_1$, $\gamma_2$, $l_*$ as in Corollary~\ref{frc:cor:transposition}. Choose pairwise
disjoint clusters in increasing order: a cluster $C_\gamma$ for each slot $\gamma$ with $\tau(\gamma)<\ell$, ordered as
$\tau$ prescribes; then clusters $Y_1<\dots<Y_m$; then a cluster $C_\gamma$ for each slot $\gamma$ with
$\tau(\gamma)>\ell+1$, ordered as $\tau$ prescribes. Let
\begin{align*}
 \cC_{\gamma_1}&=\{Y_i:f_i\in J_1\},& \mathrm{lev}^o_{\gamma_1}(Y_i)&=\ord a(f_i),& \mathrm{lev}^w_{\gamma_1}(Y_i)&=\ord x(f_i),\\
 \cC_{\gamma_2}&=\{Y_i:f_i\in J_2\},& \mathrm{lev}^o_{\gamma_2}(Y_i)&=\ord b(f_i),& \mathrm{lev}^w_{\gamma_2}(Y_i)&=\ord y(f_i),
\end{align*}
and $\cC_\gamma=\{C_\gamma\}$ with both levels $0$ for $\gamma\ne\gamma_1,\gamma_2$; these families are pairwise disjoint
since $J_1\cap J_2=\emptyset$. Let $g_\gamma$ and $h_\gamma$ be as in Lemma~\ref{frc:lem:clusters}, and
$\zeta=\widetilde Q(g_1,\dots,g_k)(h_1\wedge\dots\wedge h_k)\in\widehat\Lambda^k_\pi D$.

\emph{Step 2 (a bounded slice).} As $\vartheta$ has order $1$, (A) gives $g_{\gamma_1}=g'-g''$ with
$g'=\sum_{f_i\in J_1}a(f_i)e_{Y_i(1)}$ and $g''$ the same sum with $Y_i(3)$ in place of $Y_i(1)$, two relabelings of $a$;
so $\norm{\mu_{g_{\gamma_1}}}_{E\to D}\le2\norm{\mu_a}_{E\to D}$ by Lemma~\ref{obs:lem:multiplier}(3). Likewise
$\norm{\mu_{g_{\gamma_2}}}_{E\to D}\le2\norm{\mu_b}_{E\to D}$ and $\norm{\mu_{g_\gamma}}_{E\to D}\le2$ for the other slots,
and since $h_\gamma$ is a combination, with the coefficients $w_1,\dots,w_4$, of four relabelings of $x$, of $y$ or of a
unit vector,
$\norm{h_{\gamma_1}}_E\le4\norm x_E$, $\norm{h_{\gamma_2}}_E\le4\norm y_E$ and $\norm{h_\gamma}_E\le4$. So
\eqref{frc:eq:mult-bound} gives $\norm\zeta\le64\cdot8^{k-2}C_QN_0$. Let $\Upsilon\colon\widehat\Lambda^k_\pi D\to
D\widehat\otimes_\pi D$ be the slice of Lemma~\ref{fib:lem:slice} for the slots $\gamma_1,\gamma_2$ and
$c_\gamma=C_\gamma(1)$, let $R=\{Y_i(1):f_i\in J_1\}$ and $R'=\{Y_i(1):f_i\in J_2\}$, and let $z\in\Kh^{R\times R'}$ be the
restriction of $\Upsilon\zeta$ to $R\times R'$ (fact (P3)). Then, in the completion of $D_R\otimes_\pi D_{R'}$,
\begin{equation}\label{frc:eq:slice}
 \norm z\le64\cdot8^{k-2}k(k-1)\,C_QN_0 .
\end{equation}

\emph{Step 3 (reading all coefficients).} Let $t=Y_i(1)\in R$ and $t'=Y_{i'}(1)\in R'$. By Lemma~\ref{fib:lem:slice},
$z(t,t')=\Omega_\zeta(c)$, where $c$ has $Y_i(1)$ in slot $\gamma_1$, $Y_{i'}(1)$ in slot $\gamma_2$ and $C_\gamma(1)$ in
every other slot $\gamma$. The clusters $C_\gamma$ lie below or above all of $Y_1,\dots,Y_m$ as $\tau$ prescribes, so the
output clusters have order pattern $\sigma=\tau$ if $f_i<f_{i'}$ and $\sigma=\tau'$ if $f_i>f_{i'}$. The level sum of
Lemma~\ref{frc:lem:clusters} is $\lambda=\ord a(f_i)+\ord x(f_i)+\ord b(f_{i'})+\ord y(f_{i'})$, and
$\vartheta^\lambda=(ax)(f_i)\,(by)(f_{i'})$ by~(A). For $\sigma'\in S_k$ put
$\hat\chi(\sigma')=\sum_{l=-l_0}^{l_1}\chi_l(\sigma')X^l\in\kappa((X))$. For $l\le l_1$,
$\coeff_l(\vartheta^\lambda\hat\chi(\sigma))=\sum_{l'\ge0}\coeff_{l'}(\vartheta^\lambda)\chi_{l-l'}(\sigma)$, since
$\vartheta^\lambda$ has nonnegative order and $\chi_{l''}=0$ for $l''<-l_0$. So by Lemma~\ref{frc:lem:clusters},
$z(t,t')-\vartheta^\lambda\hat\chi(\sigma)$ has no nonzero coefficient of index at most $l_1$, and
\[
 \abs{z(t,t')-(ax)(f_i)\,(by)(f_{i'})\,\hat\chi(\sigma)}\le\varrho^{l_1+1}.
\]
The relabeling $Y_i(1)\mapsto f_i$ of rows and columns is increasing on $R\cup R'$, and by (P2) and continuity it is an
isometry onto the completion of $D_{J_1}\otimes_\pi D_{J_2}$; let $z^\sharp$ be the image of $z$. With $\tilde a=ax$ and
$\tilde b=by$, as matrices indexed by $J_1\times J_2$,
\begin{equation}\label{frc:eq:reading}
 z^\sharp=\hat\chi(\tau')\,(\tilde a\otimes\tilde b)+\hat\delta\cdot\cP_<(\tilde a\otimes\tilde b)+z_1,\qquad
 \hat\delta=\hat\chi(\tau)-\hat\chi(\tau')=\sum_{l=-l_0}^{l_1}\delta_lX^l,
\end{equation}
where $z_1\in\Kh^{J_1\times J_2}$ has entries of absolute value at most $\varrho^{l_1+1}$. The lowest nonzero coefficient
of $\hat\delta$ is $\delta_{l_*}$ with $l_*\le0$, so $\abs{\hat\delta}=\varrho^{l_*}\ge1$, and $\abs{\hat\chi(\tau')}\le\varrho^{-l_0}$.

\emph{Step 4 (conclusion under~(A)).} In the completion of $D_{J_1}\otimes_\pi D_{J_2}$, a $\Kh$\nobreakdash-normed space
by (P3), $\norm{\tilde a\otimes\tilde b}\le\norm{ax}_D\norm{by}_D\le N_0$ by Lemma~\ref{obs:lem:multiplier}(2),
$\norm{z_1}\le m^2\varrho^{l_1+1}\le N_0$ by (P3) and the choice of $l_1$, and $\norm{z^\sharp}=\norm z$ is bounded
by~\eqref{frc:eq:slice}. By~\eqref{frc:eq:reading},
\[
 \norm{\cP_<(\tilde a\otimes\tilde b)}=\abs{\hat\delta}^{-1}\norm{z^\sharp-\hat\chi(\tau')(\tilde a\otimes\tilde b)-z_1}
 \le\bigl(64\cdot8^{k-2}k(k-1)\,C_Q+\varrho^{-l_0}+1\bigr)N_0=C_\star N_0 .
\]
By (P1), the left side is $\norm{\cP_<((ax)\otimes(by))}_{D\otimes_\pi D}$. The constant does not depend on $I$, $\tau$,
$\ell$ or $l_1$, which were chosen after $a,x,b,y$.

\emph{Step 5 (removing~(A)).} Both sides of~\eqref{obs:eq:ptb} are homogeneous of degree one in each of $a,x,b,y$. The
changes (a) and (b) below do not increase $\norm{\mu_a}_{E\to D}$, $\norm x_E$, $\norm{\mu_b}_{E\to D}$ or $\norm y_E$
(Lemma~\ref{obs:lem:multiplier}(1),(3) and symmetry of the norms), and they keep $\cP_<((ax)\otimes(by))$ or change it by an
isometry of (P2).
(a)~\emph{Supports.} Replace $(a,x)$ by $(a1_J,x1_J)$ with $J=\operatorname{supp}(ax)$, and likewise $(b,y)$; then
$\operatorname{supp}a=\operatorname{supp}x=J_1$ and $\operatorname{supp}b=\operatorname{supp}y=J_2$. If
$J_1\cap J_2\ne\emptyset$, list $J_1\cup J_2$ as $f_1<\dots<f_m$ and move the entries of $a$ and $x$ at $f_i\in J_1$ to
the position $2i$ and those of $b$ and $y$ at $f_i\in J_2$ to the position $2i-1$; since $2i<2i'-1$ exactly when
$i<i'$, this relabels the rows and columns of $\cP_<((ax)\otimes(by))$, and the new supports are disjoint.
(b)~\emph{Units.} Since $\abs{K^\times}=\abs\vartheta^\Z$, each nonzero entry of $a$ is
$a(t)=\theta_a(t)\vartheta^{\ord a(t)}$ with $\abs{\theta_a(t)}=1$. Replace $a$ by the vector $a'$ with the entries
$\vartheta^{\ord a(t)}$ on $\operatorname{supp}a$ and $0$ elsewhere, and likewise $x,b,y$; this multiplies the rows of
$\cP_<((ax)\otimes(by))$ by $(\theta_a\theta_x)^{-1}$ and its columns by $(\theta_b\theta_y)^{-1}$.
(c)~\emph{Levels.} By homogeneity we may multiply $a,x,b,y$ by powers of $\vartheta$ that make all exponents
nonnegative.
After (a)--(c), (A) holds, and Steps~1--4 give~\eqref{obs:eq:ptb} with constant $C_\star$.
\end{proof}

\begin{proof}[Proof of Lemma~\ref{obs:lem:single-vector}]
We may assume $x\ne0$. Let $R=\operatorname{supp}x$, and let $z\in K^{R\times R}$ have the entries $z(t,t')=x(t)x(t')$
for $t<t'$ and $0$ otherwise; $\norm\cdot_\pi$ is the norm of Section~\ref{frc:ssec:trace}.

\emph{Step 1 (reduction to matrices).} Restricting to $R$ the vectors of a representation
$\cP_<(x\otimes x)=\sum_ag_a\otimes h_a$ with $g_a,h_a\in D$ gives a representation
$z=\sum_a(g_a|_R)(h_a|_R)^{\mathsf T}$ with $\norm{g_a|_R}_\infty\le\norm{g_a}_D$ and $\norm{h_a|_R}_\infty\le\norm{h_a}_D$.
Hence $\norm z_\pi\le\norm{\cP_<(x\otimes x)}_{D\otimes_\pi D}$.

\emph{Step 2 (normalization).} Write $\norm x_\infty=\varrho^c$ with $c\in\Z$ and put $x'=\vartheta^{-c}x$; the matrix
$z'=\vartheta^{-2c}z$ is formed from $x'$ in the same way, and $\norm z_\pi=\varrho^{2c}\norm{z'}_\pi$. Since
$\abs{K^\times}=\abs\vartheta^\Z$, each nonzero entry is $x'(t)=\theta_t\vartheta^{\lambda_t}$ with an integer $\lambda_t\ge0$
and $\abs{\theta_t}=1$, so
$z'\in\mathcal O^{R\times R}$. For $i\ge0$ let $m_i$ be the number of $t\in R$ with $\lambda_t=i$.

\emph{Step 3 (one level).} Fix $i$ with $m_i\ge2$, and let $t_1<\dots<t_{m_i}$ be the $t\in R$ with $\lambda_t=i$. The
submatrix $z'_i$ of $z'$ with rows $t_1,\dots,t_{m_i-1}$ and columns $t_2,\dots,t_{m_i}$ has the entry
$\theta_{t_a}\theta_{t_{b+1}}\vartheta^{2i}$ in position $(a,b)$ if $a\le b$, and $0$ otherwise. As $\vartheta=X\eta_0$
with a unit $\eta_0$ of $\mathcal O$, $z'_i=X^{2i}\eta_0^{2i}\Theta_1U\Theta_2$, where $\Theta_1$ and $\Theta_2$ are
diagonal with the units $\theta_{t_a}$ and $\theta_{t_{b+1}}$ and $U$ is the upper unitriangular
$(m_i-1)\times(m_i-1)$ matrix of ones; the reduction of $\eta_0^{2i}\Theta_1U\Theta_2$ modulo $X$ is invertible.
Lemma~\ref{frc:lem:trace-norm}(4),(3) gives $\rank_{2i}(z')\ge\rank_{2i}(z'_i)=m_i-1$.

\emph{Step 4 (conclusion).} The indices $2i$ are distinct, so the level formula of Lemma~\ref{frc:lem:trace-norm}(2) gives
$\norm{z'}_\pi\ge(1-\varrho)\sum_{i:m_i\ge1}\varrho^{2i}(m_i-1)$, and
\[
 \norm{x'}_2^2=\sum_{i:m_i\ge1}(m_i-1)\varrho^{2i}+\sum_{i:m_i\ge1}\varrho^{2i}
 \le\frac{\norm{z'}_\pi}{1-\varrho}+\frac1{1-\varrho^2}.
\]
Multiplying by $\varrho^{2c}=\norm x_\infty^2$ and using Step~1 gives the lemma.
\end{proof}

\begin{proof}[Proof of Proposition~\ref{obs:prop:l2-domination}]
For $m\ge1$ and $1\le r'\le\infty$ put
\[
 \psi_{r'}(m)=\sup\{\norm\xi_{r'}/\norm\xi_E:\xi\in c_{00}(\N,K)\setminus\{0\},\ \abs{\operatorname{supp}\xi}\le m\},
\]
which is finite, since $\norm\xi_{r'}\le m^{1/r'}\norm\xi_\infty\le m^{1/r'}\norm\xi_E$ by
Lemma~\ref{obs:lem:multiplier}(1); and $\psi_\infty(m)\le1$.

\emph{Step 1 (multipliers of indicators).} For finite $J\subseteq\N$ and $x\in E$, $\mu_{1_J}x=x1_J$ and
$\norm{x1_J}_D\le C'\norm{x1_J}_r\le C'\psi_r(\abs J)\norm{x1_J}_E\le C'\psi_r(\abs J)\norm x_E$, so
$\norm{\mu_{1_J}}_{E\to D}\le C'\psi_r(\abs J)$.

\emph{Step 2 (one inequality for each support size).} Let $x\in c_{00}(\N,K)\setminus\{0\}$ with $J=\operatorname{supp}x$
of size at most $m$. \textup{(PTB)} with $a=b=1_J$ and $y=x$, and Step~1, give
$\norm{\cP_<(x\otimes x)}_{D\otimes_\pi D}\le CC'^2\psi_r(m)^2\norm x_E^2$; Lemma~\ref{obs:lem:single-vector} and
$\norm x_\infty\le\norm x_E$ give $\norm x_2^2\le(\mathfrak a\psi_r(m)^2+\mathfrak b)\norm x_E^2$. So
\begin{equation}\label{frc:eq:support-size}
 \psi_2(m)^2\le\mathfrak a\,\psi_r(m)^2+\mathfrak b\qquad(m\ge1).
\end{equation}

\emph{Step 3 (interpolation).} If $r=\infty$, then $\psi_\infty\le1$ and~\eqref{frc:eq:support-size} give
$\psi_2(m)^2\le\mathfrak a+\mathfrak b=C_2^2$. Let $r<\infty$. For $\xi$ with $\abs{\operatorname{supp}\xi}\le m$,
$\norm\xi_r^r\le\norm\xi_\infty^{r-2}\norm\xi_2^2\le\psi_2(m)^2\norm\xi_E^r$, so $\psi_r(m)^2\le\psi_2(m)^{4/r}$
and~\eqref{frc:eq:support-size} gives $\psi_2(m)^2\le\mathfrak a\,\psi_2(m)^{4/r}+\mathfrak b$. If
$\psi_2(m)^2>2\mathfrak b$, then $\psi_2(m)^2<2\mathfrak a\,\psi_2(m)^{4/r}$, so $\psi_2(m)^2<(2\mathfrak a)^{r/(r-2)}$.
In both cases $\psi_2(m)^2\le C_2^2$.

\emph{Step 4 (all of $E$).} So $\norm x_2\le C_2\norm x_E$ for $x\in c_{00}(\N,K)$, and for $x\in E$,
$\norm x_2=\sup_n\norm{x1_{[n]}}_2\le C_2\sup_n\norm{x1_{[n]}}_E=C_2\norm x_E$ by
Definition~\ref{obs:def:sequence-spaces}.
\end{proof}

\begin{proof}[Proof of Proposition~\ref{obs:prop:gamma-ptb}]
Write $c_0=c_0(\N,K)$, a symmetric sequence space, and let $g_k(E)$, finite by $(\Gamma^k_E)$, be the infimum of the norms
of projections of $\widehat\otimes^k_\pi E$ onto $\widehat\Gamma^k_\pi E$ (Definition~\ref{ob:def:sequences}).

\emph{Step 1 (a multiplier lift for $(c_0,E)$).} Let $Z=\cL(\widehat\otimes^k_\pi c_0,\widehat\Lambda^k_\pi E)$ and let
$W\subseteq Z$ be the closed subspace of operators vanishing on $\cN^k_{c_0}=\ker\widehat\wedge^k_{c_0}$. Since
$\norm{\mu_a}_{c_0\to E}=\norm a_E$ (Lemma~\ref{obs:lem:multiplier}(4)), $c_{00}(\N,K)$ is dense in $E$ and $Z$ is a Banach space, the map
$\Phi(a_1,\dots,a_k)=\widehat\wedge^k_E\circ(\mu_{a_1}\otimes\dots\otimes\mu_{a_k})$ extends from $c_{00}(\N,K)^k$ to a
$k$\nobreakdash-linear $\Phi\colon E^k\to Z$ of norm at most $1$. Its diagonal lies in $W$, since
$\Phi(a,\dots,a)=\widehat\Lambda^k_\pi\mu_a\circ\widehat\wedge^k_{c_0}$ for $a\in c_{00}(\N,K)$ by naturality of the wedge
and $W$ is closed. Lemma~\ref{ob:lem:parameter}(3), with $P=E$, $n=k$ and $\delta=1$, gives a bounded
$k$\nobreakdash-linear $\Phi'\colon E^k\to W$ with the diagonal of $\Phi$ and $\norm{\Phi'}\le g_k(E)+1$. By
Lemma~\ref{ob:lem:metric-surjection}, $T\mapsto T\circ\widehat\wedge^k_{c_0}$ is a linear isometry of
$\cL(\widehat\Lambda^k_\pi c_0,\widehat\Lambda^k_\pi E)$ onto $W$. For $a_1,\dots,a_k\in c_{00}(\N,K)$ let
$\widetilde Q(a_1,\dots,a_k)$ be the restriction to $\Lambda^k_\pi c_0$ of the preimage of $\Phi'(a_1,\dots,a_k)$. Then
$\widetilde Q$ is $k$\nobreakdash-linear, $\widetilde Q(a,\dots,a)=\Lambda^k_\pi\mu_a$, and
\[
 \norm{\widetilde Q(a_1,\dots,a_k)}\le(g_k(E)+1)\prod_{\gamma=1}^k\norm{\mu_{a_\gamma}}_{c_0\to E}.
\]

\emph{Step 2 (the triangle bound).} By Lemma~\ref{obs:lem:multiplier}(1),(3), $\widetilde Q$ is a multiplier lift with
$G=c_0$, target $E$ and $l_0=\max(0,\lceil\log_{1/\varrho}(g_k(E)+1)\rceil)$, and it satisfies~\eqref{frc:eq:mult-bound}
for $(c_0,E)$ with $C_Q=g_k(E)+1$. So Steps~1--5 of the proof of Theorem~\ref{obs:thm:ptb} give \textup{(PTB)} for
$(c_0,E)$, and the choice for $G=c_0$ in Step~1 of the proof of Theorem~\ref{cls:thm:symmetric} bounds $\cP_<$ on
$E\otimes_\pi E$.
\end{proof}

\subsection{Nonarchimedean fibers}\label{frc:ssec:nonarch}
\begin{proof}[Proof of Lemma~\ref{fib:lem:marking}]
\emph{Step 1 (a normalized basis).} Extend an equivalent nonarchimedean norm of $E$ to $\widehat E$ and round it up to
$\varrho^\Z\cup\{0\}$, as in the proof of Corollary~\ref{cor:discrete-targets}; this gives $\norm\cdot'$. We construct
$f_1,f_2,\dots\in\widehat E$ with $\norm{\sum_n\xi_nf_n}'=\max_n\abs{\xi_n}$. Given $f_1,\dots,f_n$, their span $V_n$ is
closed, being finite-dimensional over the complete field $\Kh$ (Lemma~\ref{lem:findim}). For $v\in\widehat E\setminus V_n$
the distances $\norm{v-w}'$, $w\in V_n$, lie in $\varrho^\Z$ and are at least $\operatorname{dist}(v,V_n)>0$, so some
$w_0$ attains their minimum, and the nonarchimedean inequality gives
$\norm{\lambda(v-w_0)+w}'=\max(\abs\lambda\norm{v-w_0}',\norm w')$ for $\lambda\in\Kh$ and $w\in V_n$. Let $f_{n+1}$ be a
multiple of $v-w_0$ of norm one.

\emph{Step 2 (the projection).} The closed span $V_\infty$ of the $f_n$ is isometric to $c_0(\N,\Kh)$, complete with
distances in $\varrho^\Z\cup\{0\}$, hence spherically complete (Lemma~\ref{lem:discdist}); the extension argument in the
proof of Theorem~\ref{thm:spherical}, over $\Kh$, gives a contractive projection $\Pi$ of $\widehat E$ onto $V_\infty$.
Choose $e_n\in E$ with $\norm{e_n-f_n}'\le\varrho$. Then $\norm{\sum_n\xi_n(e_n-f_n)}'<\max_n\abs{\xi_n}$ for $\xi\ne0$,
which gives the norm formula for the $e_n$. The operator $\Pi'$ on $V_\infty$ with
$\Pi'(\sum_n\xi_nf_n)=\Pi(\sum_n\xi_ne_n)$ satisfies $\norm{\Pi'-\id_{V_\infty}}\le\varrho<1$, so it is invertible with
$\norm{\Pi'^{-1}}\le1$. Let $\theta_E(x)$ be the coordinate sequence of $\Pi'^{-1}\Pi x$ in the basis $(f_n)$; then
$\theta_E(e_n)$ is the $n$th unit vector. If $E$ is complete, then $E=\widehat E$ is a $\Kh$\nobreakdash-Banach space and
the series defining $\widehat\psi_E$ converge; if $K$ is complete, $K=\Kh$.
\end{proof}

\begin{proof}[Proof of Theorem~\ref{fib:thm:marking}]
\emph{Step 1 (the form $\Phi$).} Let $F_1$ and $\eta$ be as in Appendix~\ref{app:laurent} for $K_1=\Kh$, so that
$E_1=\linf_{\Kh}$, and let $\eta'$ be $\eta$ composed with the map $\widehat\Lambda^k_\pi\linf_{\Kh}\to F_1$ induced by
the identity of $\linf_{\Kh}$, a contraction since a decomposition over $K$ is one over $\Kh$ with the same sum
$\sum_j\prod_i\norm{x_{ji}}$. Suppose $Q$ exists, and for $a,x\in c_{00}(\N,K)^k$ put
\[
 \Phi(a;x)=\theta_F\bigl(Q(a_1,\dots,a_k)(\psi x_1,\dots,\psi x_k)\bigr)\in\widehat\Lambda^k_\pi\linf_{\Kh}.
\]
It is $k$\nobreakdash-linear in $a$ and in $x$, alternating in $x$, of norm at most $\norm{\theta_F}\norm Q\norm\psi^k$,
and $\Phi(a,\dots,a;x)=\omega_{\linf_{\Kh}}(ax_1,\dots,ax_k)$ by~\eqref{fib:eq:marking}.

\emph{Step 2 (reduction to $\kappa$).} For $a,x$ with entries in $\F_p$, $\Psi(a;x)=\eta'(\Phi(a;x))\in
\Lambda^k_\kappa(\kappa^\N)$ is $\F_p$\nobreakdash-multilinear and alternating in $x$, and by Proposition~\ref{prop:supp},
$\sdim\Psi(a;x)\le\lfloor kM_\varrho\norm{\theta_F}\norm Q\norm\psi^k\rfloor$ when all entries lie in $\{0,\pm1\}$. As in
Step~1 of the proof of Theorem~\ref{fib:thm:ell2}, polarizing the diagonal identity over the infinite field $K$ and
applying $\eta'$, which sends $\omega_{\linf_{\Kh}}(y_1,\dots,y_k)$ to $y_1\wedge\dots\wedge y_k$ when every $y_j$ has
entries in $\kappa$ (proof of Lemma~\ref{lem:Psi}), gives~\eqref{eq:Pol1} for $\Psi$ on such inputs. Let $\Psi^\kappa$ be
the $\kappa$\nobreakdash-multilinear map on $(\kappa^{(\N)})^{2k}$ that agrees with $\Psi$ on unit vectors. Antisymmetry in
the last $k$ slots and~\eqref{eq:Pol1} are identities between $\kappa$\nobreakdash-multilinear expressions that hold on
unit vectors, so they hold for $\Psi^\kappa$, and $\Psi^\kappa=\Psi$ on inputs with entries in $\{0,\pm1\}$. Since $k!=0$
in $\kappa$, this contradicts Theorem~\ref{thm:FF}(1).

\emph{Step 3 (consequences).} If $A^k_{E,E';F}$ were analytic at a point, Proposition~\ref{prop:homogeneous-lift} would
give a bounded lift $Q_0$, and $Q(a_1,\dots,a_k)=Q_0(u_{a_1},\dots,u_{a_k})(\nu)$ would be excluded. The last assertion
follows from Lemma~\ref{fib:lem:test}.
\end{proof}

\begin{proof}[Proof of Theorem~\ref{fib:thm:local}]
\emph{Step 1 (finite dimension).} If $\dim_{\Kh}\widehat E<\infty$, then $\widehat E$ has a $\Kh$\nobreakdash-basis
$(e_i)$ with continuous coordinates $(\varepsilon_i)$ (Lemma~\ref{lem:findim}). For $u_1,\dots,u_k\in\cL(E,E')$, with
extensions $\widehat u_i$ to the completions, put
\[
 Q(u_1,\dots,u_k)(x_1\wedge\dots\wedge x_k)=\sum_{i_1<\dots<i_k}\det\bigl(\varepsilon_{i_a}(x_b)\bigr)_{a,b}\,
 \widehat\omega_{E'}(\widehat u_1e_{i_1},\dots,\widehat u_ke_{i_k}),
\]
where the continuous extension $\widehat\omega_{E'}$ of $\omega_{E'}$ is $\Kh$\nobreakdash-multilinear by density.
Expanding $x_b=\sum_i\varepsilon_i(x_b)e_i$ shows, as in Proposition~\ref{prop:split-lift}(2), that $Q$ is a bounded lift
of $u\mapsto\Lambda^k_\pi u$ with values in $\cL(\Lambda^k_\pi E,\widehat\Lambda^k_\pi E')$. For $u_i\in\cL(E',E)$ use the
terms $\det(\varepsilon_{i_a}(u_ax_b))_{a,b}\,\widehat\omega_E(e_{i_1},\dots,e_{i_k})$ instead. If $E$ is complete, then
$e_i\in E$, and moving the scalar into the first argument shows that these lifts take values in the uncompleted spaces.
This proves the positive directions; in~(3) with $\dim E'<\infty$, exchange the roles of $E$ and $E'$ in the second
formula.

\emph{Step 2 (infinite dimension).} Let $\dim_{\Kh}\widehat E=\infty$. Take $e_n$, $\theta_E$ and, for complete $E$,
$\widehat\psi_E$ from Lemma~\ref{fib:lem:marking}, and put $\psi_E(\xi)=\sum_n\xi_ne_n$ on $c_{00}(\N,K)$.
Theorem~\ref{fib:thm:marking} applies to the following data, which satisfy~\eqref{fib:eq:marking}:
\begin{enumerate}
\item[(a)] the pair $(E,c_0(\N,\Kh))$, with $\psi=\psi_E$, $\theta$ the inclusion, and $u_a(x)=a\,\theta_E(x)$;
\item[(b)] the pair $(c_0(\N,K),E)$, with $\psi$ the inclusion, $\theta=\theta_E$, and $u_a(x)=\psi_E(ax)$;
\item[(c)] for complete $E$, the pair $(E,E)$, with $\psi=\psi_E$, $\theta=\theta_E$, and
$u_a(x)=\widehat\psi_E(a\,\theta_E(x))$;
\item[(d)] for complete $K$ and infinite-dimensional $E,E'$, the pair $(E,E')$, with $\psi=\psi_E$,
$\theta=\theta_{E'}$, and $u_a(x)=\psi_{E'}(a\,\theta_E(x))$, since $\theta_E$ then takes values in $c_0(\N,K)$.
\end{enumerate}
Over complete $K$, $\dim_{\Kh}\widehat E=\dim E$, since finite-dimensional spaces are complete. An equivalent
nonarchimedean norm on the values would make the action analytic by Corollary~\ref{cor:discrete-targets}, since
$\abs{K^\times}$ is discrete.
\end{proof}

\subsection{Proofs of the classification results}\label{frc:ssec:classification}
As in Section~\ref{sec:classification}, $c_{00}$, $c_0$ and $\ell^r$ stand for $c_{00}(\N,K)$, $c_0(\N,K)$ and
$\ell^r(\N,K)$.

\begin{proof}[Proof of Theorem~\ref{cls:thm:symmetric}]
\emph{Step 1 ((1)$\Rightarrow$(3) and (2)$\Rightarrow$(3)).} Let $G\in\{E,c_0\}$ with $(G,E)$ good for Banach values in
a degree $k\ge p$; Theorem~\ref{obs:thm:ptb} gives \textup{(PTB)} for $(G,E)$ with a constant $C$. For $x,y\in c_{00}$,
\eqref{obs:eq:ptb} with $a=1_{\operatorname{supp}x}$ and $b=1_{\operatorname{supp}y}$ if $G=E$, and with
$(x,1_{\operatorname{supp}x},y,1_{\operatorname{supp}y})$ in place of $(a,x,b,y)$ if $G=c_0$, gives
$\norm{\cP_<(x\otimes y)}_{E\otimes_\pi E}\le C\norm x_E\norm y_E$, since $\norm{\mu_{1_S}}_{E\to E}\le1$, $\norm{1_S}_{c_0}=1$
and $\norm{\mu_x}_{c_0\to E}=\norm x_E$ (Lemma~\ref{obs:lem:multiplier}(1),(4)).

\emph{Step 2 ((3)$\Rightarrow$(4), (5), (6), (7)).} Condition~(3) is the hypothesis of
Proposition~\ref{spl:prop:ordering} for the unit vectors, a linearly independent family indexed by $\N$ with dense span
$c_{00}$. So $(\mathrm S^n_E)$, $(\mathrm R^n_E)$ and $(\Gamma^n_E)$ hold for every $n\ge1$, which gives (4), (5) and (6) with
$n=p$, and Proposition~\ref{fib:prop:split}(1),(2) gives~(7).

\emph{Step 3 ((4)$\Rightarrow$(1), (5)$\Rightarrow$(1), (6)$\Rightarrow$(3), (7)$\Rightarrow$(1), (2)).}
Proposition~\ref{fib:prop:split}(1), respectively~(2), with $D=E$ gives the first two implications, and
Proposition~\ref{obs:prop:gamma-ptb} the third; for the last take $D=E$, respectively $D=c_0$, and the degree $p$.

\emph{Step 4 ($\ell^2$-domination).} Let $C$ be the constant in~(3). Lemma~\ref{obs:lem:single-vector} with $D=E$ and
$\norm x_\infty\le\norm x_E$ (Lemma~\ref{obs:lem:multiplier}(1)) gives $\norm x_2\le C_2\norm x_E$ for $x\in c_{00}$, with
$C_2^2=C/(1-\varrho)+1/(1-\varrho^2)$, hence for $x\in E$ by taking suprema over the $x1_{[n]}$
(Definition~\ref{obs:def:sequence-spaces}).

\emph{Step 5 (the fundamental function).} We may assume $C>0$. Let $m\ge1$ and
$z_m=\cP_<(1_{[m]}\otimes1_{[m]})=\sum_{t<t'}e_t\otimes e_{t'}$, the sum over $t,t'\in[m]$. By~(3) and fact (P1) of
Section~\ref{frc:ssec:ptb}, $z_m$ has norm at most $C\varphi_E(m)^2$ in $E_{[m]}\otimes_\pi E_{[m]}$; the identity of
$E_{[m]}$ onto $\ell^2([m],K)$ has norm at most $C_2$ by Step~4, so $z_m$ has norm at most $C_2^2C\varphi_E(m)^2$ in
$\ell^2([m],K)\otimes_\pi\ell^2([m],K)$. The matrix $\mathsf Z=\coeff_0z_m$ is the strictly upper triangular matrix of
ones, so every off-diagonal entry of $\mathsf Z-\mathsf Z^{\mathsf T}$ is $\pm1$, and Lemmas~\ref{fib:lem:tournament}
and~\ref{fib:lem:coeff}(1) give
\[
 \frac{m\ln m}{1+\ln(2\abs\kappa)}\le\operatorname{cost}(\mathsf Z)\le\frac{C_2^2C\varphi_E(m)^2}{1-\varrho^2},
\]
so $\varphi_E(m)^2\ge c\,m\ln m$ with $c>0$ independent of $m$.
\end{proof}

\begin{proof}[Proof of Theorem~\ref{cls:thm:rnorm}]
(1)$\Rightarrow$(2) is Theorem~\ref{obs:thm:ptb}, and (4)$\Rightarrow$(1) follows with $k=p$.

(2)$\Rightarrow$(3). Let $\nu$ be an equivalent $r$\nobreakdash-norm on $D$, or an equivalent nonarchimedean norm, in
which case put $r=\infty$, scaled so that $\nu(y)\le\norm y_D\le C''\nu(y)$ for $y\in D$. Since $\nu(e_t)\le1$, iterating
$\nu(y+y')^r\le\nu(y)^r+\nu(y')^r$, respectively $\nu(y+y')\le\max(\nu(y),\nu(y'))$, gives
$\norm\xi_D\le C''\nu(\xi)\le C''\norm\xi_r$ for $\xi\in c_{00}$. As $\norm y_\infty\le\norm y_D$
(Lemma~\ref{obs:lem:multiplier}(1)), Proposition~\ref{obs:prop:l2-domination} with $C'=C''$ gives~(3).

(3)$\Rightarrow$(4). The unit vectors of $E$ are a linearly independent family indexed by $\N$, of norm one with dense
span $c_{00}$, and $\ell^2$\nobreakdash-dominated by~(3) (Definition~\ref{obs:def:sequence-spaces}). As $K$ is
nonarchimedean, Theorem~\ref{spl:thm:relative-ordering}(2) for $r<\infty$, and
Corollary~\ref{spl:cor:relative-ordering-na} for a nonarchimedean norm, give~(4).
\end{proof}

\begin{proof}[Proof of Corollary~\ref{cls:cor:lr-table}]
\emph{Positive cases.} If $r_i<2$ for some $i$, then $\ell^{r_i}$ has a lower $r_i$\nobreakdash-estimate with its unit
vectors and $C=1$ (the examples after Theorem~\ref{fib:thm:lower}), and Theorem~\ref{fib:thm:lower}(3) makes the pair good
for Banach values in every degree. If $r_1=2<r_2$, the unit vectors of $\ell^2$ are $\ell^2$\nobreakdash-dominated, of norm one
with dense span; $\ell^{r_2}$ is $r_2$\nobreakdash-normed by its own norm if $r_2<\infty$, as noted
after Corollary~\ref{spl:cor:relative-ordering-na}, and $c_0$ has a nonarchimedean norm, so
Theorem~\ref{spl:thm:relative-ordering}(2), respectively Corollary~\ref{spl:cor:relative-ordering-na}, makes the pair good
for Banach values in every degree.

\emph{Negative cases.} For the remaining pairs, $r_2=2\le r_1$ or $r_1,r_2>2$, Corollary~\ref{fib:cor:pairs}
excludes every degree $k\ge p$ and Corollary~\ref{fib:cor:degrees}(1) gives every degree $k<p$. So
$\fD_B(\ell^{r_1},\ell^{r_2})$ is $\{1,2,\dots\}$ in the positive cases and $\{1,\dots,p-1\}$ otherwise, which gives the
equivalence and the table.
\end{proof}

\subsection{Spreading spaces}\label{frc:ssec:spreading}
For a finite set $S\subseteq\N$, a map $\psi\colon S\to\N$ that is increasing and $x\in c_{00}(\N,K)$ supported in $S$,
put $\psi_*x=\sum_{t\in S}x(t)e_{\psi(t)}$. If $E$ is spreading, then $\abs{\psi_*x}=\sum_{t\in S}\abs{x(t)}e_{\psi(t)}$
and Definition~\ref{cls:def:spreading} give $\norm{\psi_*x}_E=\norm x_E$; in particular $\norm{1_S}_E=\varphi_E(\abs S)$.

\begin{proof}[Proof of Theorem~\ref{cls:thm:spreading}]
\emph{Step 1 (item (1)).} In the proof of Lemma~\ref{obs:lem:multiplier}, invariance under permutations is used only
in~(3), to see that a vector $y$ supported in $\operatorname{supp}a$ and its relabeling $y^\gamma$ have the same norm.
For $\gamma$ increasing on $\operatorname{supp}a$ this is the identity $\norm{\psi_*y}=\norm y$ above, with $\psi$ the
restriction of $\gamma$ to $\operatorname{supp}a$. The rest of the proof uses only monotonicity, $\norm{e_t}=1$ and
$\norm x=\sup_n\norm{x1_{[n]}}$.

\emph{Step 2 (the facts (P1)--(P3)).} They hold for spreading $D$, with (P2) for relabelings of the rows along a map
increasing on the row support and of the columns along a map increasing on the column support: such relabelings map
$D_J$ and $D_{J'}$ isometrically onto $D_{\psi(J)}$ and $D_{\psi'(J')}$, so their tensor product is an isometry for the
projective norms, and (P1) applies.

\emph{Step 3 (Theorem~\ref{obs:thm:ptb}).} Lemmas~\ref{frc:lem:homogeneity}, \ref{frc:lem:clusters}
and~\ref{frc:lem:diagonal} and Corollary~\ref{frc:cor:transposition} concern sequence spaces $G$ and $D$ with
$\norm{e_t}_G\le1$ and $\norm\cdot_\infty\le\norm\cdot_D$, and these properties hold for spreading spaces by Step~1. In
the proof of Theorem~\ref{obs:thm:ptb}, the norms of $E$ and $D$ enter through Lemma~\ref{obs:lem:multiplier}, the facts
(P1)--(P3) and relabelings, and each relabeling is increasing on the support it is applied to: in Step~2, $f_i\mapsto
Y_i(c)$ on $J_1$ and on $J_2$, $c=1,\dots,4$, since $Y_1<\dots<Y_m$; in Step~3, $Y_i(1)\mapsto f_i$ on $R$ and on $R'$; in
Step~5(a), $f_i\mapsto2i$ on $J_1$ and $f_i\mapsto2i-1$ on $J_2$. Steps~5(b) and~(c) use only $\norm x=\norm{\abs x}$ and
homogeneity. So Steps~1 and~2 give the theorem for spreading $E$ and $D$.

\emph{Step 4 (the other results of item (2)).} The proofs of Propositions~\ref{obs:prop:l2-domination}
and~\ref{obs:prop:gamma-ptb} and of Theorems~\ref{cls:thm:symmetric} and~\ref{cls:thm:rnorm} use the symmetry of the
spaces only through Lemma~\ref{obs:lem:multiplier}, Theorem~\ref{obs:thm:ptb} and fact~(P1), together with
$\norm{e_t}_D=1$ in Theorem~\ref{cls:thm:rnorm}; Lemma~\ref{obs:lem:single-vector} concerns every sequence space $D$ with
$\norm\cdot_\infty\le\norm\cdot_D$, and $c_0$ is spreading. Corollaries~\ref{cls:cor:symmetric-examples}
and~\ref{cls:cor:dichotomy}(1),(2) follow from these results as before. Remark~\ref{cls:rem:mechanisms} rests, for these
pairs, on Proposition~\ref{spl:prop:ordering} (Step~2 of the proof of Theorem~\ref{cls:thm:symmetric}) and on the
relative ordering splitting described after Theorem~\ref{cls:thm:rnorm} (Theorem~\ref{spl:thm:relative-ordering}(1) and
Corollary~\ref{spl:cor:relative-ordering-na}), which concern arbitrary normed spaces.

\emph{Step 5 (item (3)).} Here $K$ is any nontrivially normed field of characteristic~$2$; condition~\textup{(L)} is not
used. Steps~7 and~8 of the proof of Theorem~\ref{obs:thm:char-two} use of $E$ only that it is a sequence
space with $\norm{e_t}_E=1$, $\norm x_\infty\le\norm x_E$, $\norm{\mu_a}_{E\to E}\le\norm a_\infty$ and
$\sup_nn/\bar\varphi_E(n)^2=\infty$. For spreading $E$ the first three follow from monotonicity, as there, and
$\bar\varphi_E=\varphi_E$, since $\norm{1_S}_E=\varphi_E(\abs S)$ and $\varphi_E$ is nondecreasing.
\end{proof}

\begin{proof}[Proof of the assertions of Example~\ref{cls:ex:garling}]
Write $\norm\cdot$ for $\norm\cdot_{\alpha,q}$ on $c_{00}(\N,\mathbb R)$, $\varphi$ for $\varphi_{E_{\alpha,q}}$, and put
$\mathfrak h_n=\sum_{j\le n}1/j$.

\emph{Step 1 (spreading).} Minkowski's inequality along each increasing sequence makes $\norm\cdot$ a norm. It depends
only on $\abs f$, is monotone in $\abs f$ and dominates $\norm f_\infty$, and $\norm{e_t}=1$. Let
$f=\sum_{i\le d}f_ie_{t_i}$ with $t_1<\dots<t_d$. An increasing sequence $(u_j)$ contributes only through its terms in
$\{t_1,\dots,t_d\}$, say $u_{j_s}=t_{a_s}$ with $j_1<j_2<\cdots$ and $a_1<a_2<\cdots$; as $j_s\ge s$, its sum is at most
$\sum_ss^{-\alpha}\abs{f_{a_s}}^q$, the sum of the sequence $(t_{a_s})_s$. So $\norm f^q$ is the largest value of
$\sum_ss^{-\alpha}\abs{f_{a_s}}^q$ over the increasing sequences $(a_s)$ in $\{1,\dots,d\}$, which does not depend on
$t_1,\dots,t_d$. The other conditions of Definition~\ref{cls:def:spreading} hold by the definition of $E_{\alpha,q}$.

\emph{Step 2 (decreasing profiles; item (1)).} If $\abs{g(1)}\ge\abs{g(2)}\ge\cdots$, then $\abs{g(u_j)}\le\abs{g(j)}$
for every increasing sequence $(u_j)$, so $\norm g^q=\sum_jj^{-\alpha}\abs{g(j)}^q$. In particular
$\varphi(m)^q=\sum_{j\le m}j^{-\alpha}$, which lies between $m^{1-\alpha}$ and
$1+\int_1^m\tau^{-\alpha}\,d\tau\le m^{1-\alpha}/(1-\alpha)$.

\emph{Step 3 (blocks).} For $n,M\ge1$ let $I_i=\{(i-1)n+1,\dots,in\}$ and
$x=\sum_{i\le M}\sum_{j\le n}j^{-\beta}e_{(i-1)n+j}$, a real vector. By Steps~1 and~2 each $x1_{I_i}$ has norm $\mathfrak
h_n^{1/q}$, since $j^{-\alpha}j^{-\beta q}=j^{-1}$. We show $\norm x^q\le1+\ln(Mn)+(M-1)/\alpha$. It suffices to consider
increasing sequences in the support of $x$, since other terms only raise the indices of the later ones. Fix one, let
$m_i$ be the number of its terms in $I_i$, and $\lambda_i=m_1+\dots+m_{i-1}$. The terms in $I_i$ have indices
$\lambda_i+s$, $s\le m_i$, at positions $(i-1)n+j_s$ with $j_s\ge s$, so they contribute at most
$G_i=\sum_{s\le m_i}(\lambda_i+s)^{-\alpha}s^{\alpha-1}$. If $\lambda_i=0$, then $G_i=\mathfrak h_{m_i}\le1+\ln m_i$. If
$\lambda_i\ge1$, the function $\tau\mapsto(\lambda_i+\tau)^{-\alpha}\tau^{\alpha-1}$ decreases on $(0,\infty)$, so the
substitution $\tau=\lambda_iv$ gives
\[
 G_i\le\int_0^{m_i/\lambda_i}(1+v)^{-\alpha}v^{\alpha-1}\,dv\le\frac1\alpha+\ln\frac{\lambda_i+m_i}{\lambda_i},
\]
using $(1+v)^{-\alpha}\le1$ on $(0,1]$ and $(1+v)^{-\alpha}v^{\alpha-1}\le v^{-1}$ on $[1,\infty)$. Let $i_0$ be the first
$i$ with $m_i>0$. At most $M-1$ later blocks are met, and their logarithms add up to the logarithm of the number of
terms divided by $m_{i_0}$, at most $\ln(Mn/m_{i_0})$. So the sum of the sequence is at most
$1+\ln m_{i_0}+(M-1)/\alpha+\ln(Mn/m_{i_0})$.

\emph{Step 4 (item (2): symmetric norms).} Let $x^*=x\circ\sigma$ be the decreasing rearrangement of $x$, for a
permutation $\sigma$ of $[Mn]$; it takes the value $j^{-\beta}$ at the $M$ positions $(j-1)M+1,\dots,jM$. By Step~2,
\[
 \norm{x^*}^q=\sum_{j\le n}j^{\alpha-1}\sum_{l\le M}((j-1)M+l)^{-\alpha}\ge M^{1-\alpha}\mathfrak h_n,
\]
so by Step~3 the ratio $\norm{x\circ\sigma}^q/\norm x^q$ is at least $M^{1-\alpha}\mathfrak h_n/(1+\ln(Mn)+(M-1)/\alpha)$,
which is unbounded as $n$ and then $M$ grow. For real $f\ge0$ in $c_{00}(\N,\mathbb R)$ let $f'\in c_{00}(\N,K)$ have
$f'(t)=0$ if $f(t)=0$ and $f'(t)=\vartheta^{c(t)}$ otherwise, with $c(t)$ the largest integer such that
$\abs\vartheta^{c(t)}\ge f(t)$. Then $f\le\abs{f'}<\abs\vartheta^{-1}f$, so $\norm f\le\norm{f'}\le\abs\vartheta^{-1}\norm f$,
and $(f\circ\sigma)'=f'\circ\sigma$. Hence $\norm{x'\circ\sigma}/\norm{x'}\ge\abs\vartheta\norm{x\circ\sigma}/\norm x$ is
unbounded, which a norm on $c_{00}$ invariant under permutations and equivalent to that of $E_{\alpha,q}$ would exclude.

\emph{Step 5 (item (2): lower estimates).} Let $r\in(0,2)$ and $C\ge1$, and let $x'$ be as in Step~4. The $I_i$ are
intervals, $\norm{x'1_{I_i}}\ge\norm{x1_{I_i}}=\mathfrak h_n^{1/q}$ and $\norm{x'}\le\abs\vartheta^{-1}\norm x$, so
\eqref{fib:eq:lower} would give $M^{1/r}\mathfrak h_n^{1/q}\le C\abs\vartheta^{-1}(1+\ln(Mn)+(M-1)/\alpha)^{1/q}$. For
$M=\lceil\ln n\rceil$ and large $n$ the right side is at most $C\abs\vartheta^{-1}((2+1/\alpha)\ln n)^{1/q}$ and the left
side at least $(\ln n)^{1/r+1/q}$, so $(\ln n)^{1/r}\le C\abs\vartheta^{-1}(2+1/\alpha)^{1/q}$, which fails for large $n$.

\emph{Step 6 (items (3) and (4)).} If $2\beta\le1$, then $\varphi(m)^2\le(1-\alpha)^{-2/q}m$ by~(1), and
Corollary~\ref{cls:cor:symmetric-examples}, which holds for spreading spaces by Theorem~\ref{cls:thm:spreading}(2), gives
the first assertion of~(3); if $2\beta<1$, then $\sup_mm/\varphi(m)^2=\infty$, and Theorem~\ref{cls:thm:spreading}(3)
gives the second, Steps~1 and~2 holding over every nontrivially normed field. For~(4), let $y\in c_{00}$ and let $y^*_m$
be the $m$th largest value of $\abs y$. The positions of the $m$ largest values, in increasing order, give
$\norm y^q\ge\sum_{s\le m}s^{-\alpha}(y^*_m)^q\ge m^{1-\alpha}(y^*_m)^q$, so $y^*_m\le m^{-\beta}\norm y$ and
$\norm y_2^2\le(\sum_mm^{-2\beta})\norm y^2$. For $y\in E_{\alpha,q}$ take the supremum over the $y1_{[n]}$.
\end{proof}
